% =====================================================================
%  CV — José Manuel Requena Plens — Versión ATS (ES)
%  Diseño: una columna "diseñada", ATS-friendly. Motor: LuaLaTeX + fontspec.
%  Compilar (nombre para reclutadores en el PDF de salida):
%    lualatex -jobname="Jose_Requena_Plens_Firmware_Software_2026" CV_RequenaPlensJoseManuel_SPA_ATS.tex
% =====================================================================
% Accesibilidad: PDF etiquetado (tagged) + idioma del documento.
\DocumentMetadata{lang=es-ES, pdfversion=2.0, testphase={phase-III}}
\documentclass[11pt,a4paper]{article}

% ---- Codificación / idioma ----
\usepackage{fontspec}
\usepackage[spanish,es-noindentfirst]{babel}
\usepackage{microtype}

% ---- Geometría (márgenes amplios = aire "premium") ----
\usepackage[a4paper,top=1.5cm,bottom=1.4cm,left=1.7cm,right=1.7cm]{geometry}

% ---- Color (un solo acento = morado de jmrp.io) ----
\usepackage{xcolor}
\definecolor{accent}{HTML}{B509AC}     % morado de marca (light theme)
\definecolor{ink}{HTML}{1F2328}        % texto principal (casi negro)
\definecolor{muted}{HTML}{57606A}      % texto secundario

% ---- Tipografía: Inter ----
\setmainfont{Inter}[
  UprightFont = *-Regular,
  BoldFont    = *-SemiBold,
  ItalicFont  = *-Italic,
  BoldItalicFont = *-SemiBoldItalic,
  WordSpace   = 1.15,        % ensancha el espacio entre palabras: extracción ATS más robusta
  RawFeature  = {-tlig},     % apóstrofo/comillas rectos (mejor extracción ATS)
]
\newfontfamily\headingfont{Inter}[UprightFont=*-Bold]
\color{ink}

% ---- Listas compactas ----
\usepackage{enumitem}
\setlist[itemize]{
  leftmargin=1.1em, topsep=2pt, partopsep=0pt, parsep=0pt, itemsep=2pt,
  label=\small\textcolor{accent}{\textbullet}
}

% ---- Secciones (versalitas + regla fina en acento) ----
\usepackage{titlesec}
\titleformat{\section}
  {\headingfont\large\color{accent}}
  {}{0em}{}[\vspace{-0.55em}{\color{accent}\rule{\linewidth}{0.8pt}}]
\titlespacing*{\section}{0pt}{0.6em}{0.5em}

% Evita encabezados de sección huérfanos al final de página:
% si no caben al menos ~4 líneas de contenido tras el título, salta de página.
\usepackage{needspace}
\let\jmrpsection\section
\renewcommand{\section}[1]{\needspace{4\baselineskip}\jmrpsection{#1}}

% ---- Enlaces ----
\usepackage[hidelinks]{hyperref}
\hypersetup{
  colorlinks=true, urlcolor=accent, linkcolor=accent,
  pdftitle={José Manuel Requena Plens — CV (Ingeniero de Firmware y Software)},
  pdfauthor={José Manuel Requena Plens},
  pdfsubject={Curriculum Vitae},
  pdfkeywords={firmware, embedded, C, STM32, ESP32, FreeRTOS, Go, Python, QA, CI/CD, SonarQube, Modbus, RTOS, software, DevSecOps},
  pdflang={es-ES},
}
\usepackage{fontawesome5}

% ---- Interlineado cómodo ----
\linespread{1.05}
\setlength{\parindent}{0pt}

% =====================================================================
%  Comandos de ayuda
% =====================================================================
% Entrada de experiencia/formación: {cargo}{organización}{lugar}{fechas}
\newcommand{\cventry}[4]{%
  \textbf{#1}\hfill\textbf{#4}\par
  {\color{muted}\textit{#2}\hfill\textit{#3}}\par
  \vspace{2pt}
}
% Proyecto: {nombre}{stack}{métricas}
\newcommand{\cvproject}[3]{%
  \textbf{#1}\,\textcolor{muted}{\small— #2}\hfill{\small\textcolor{accent}{#3}}\par
  \vspace{1pt}
}
% Skill: {categoría}{items}
\newcommand{\cvskill}[2]{\textbf{#1:} #2\par\vspace{2pt}}
\newcommand{\divider}{\textcolor{muted}{\ \raisebox{0.4ex}{\scriptsize$\bullet$}\ }}

\begin{document}

% =====================================================================
%  ENCABEZADO
% =====================================================================
{\headingfont\fontsize{22pt}{24pt}\selectfont José Manuel Requena Plens}\par
\vspace{2pt}
{\large\color{accent}Ingeniero de Firmware y Software}\par
\vspace{4pt}
{\small\color{muted}%
Valencia, España \divider
\href{mailto:jmrplens@gmail.com}{jmrplens@gmail.com} \divider
% TODO: teléfono (no figura en el repo) — añadir o eliminar
\href{https://jmrp.io}{jmrp.io} \divider
\href{https://github.com/jmrplens}{github.com/jmrplens} \divider
\href{https://www.linkedin.com/in/jmrplens}{linkedin.com/in/jmrplens}%
}\par
\vspace{2pt}
{\color{accent}\rule{\linewidth}{1.2pt}}

% =====================================================================
%  PERFIL
% =====================================================================
\section{Perfil}
\textbf{Ingeniero de firmware y software} con más de 3 años desarrollando firmware industrial para inversores solares (C, STM32, FreeRTOS) y una base sólida en \textbf{I+D}. Especializado en optimización de bajo nivel, comunicaciones industriales (Modbus) y automatización de pruebas sobre hardware. Mantengo proyectos \textit{open source} usados por desarrolladores de todo el mundo. Oriento mi carrera a roles de \textbf{firmware, software y QA}, aportando rigor de ingeniería y foco en la calidad.

% =====================================================================
%  EXPERIENCIA
% =====================================================================
\section{Experiencia}
\cventry{Desarrollador de Software I+D}{Power Electronics — Solar + Storage}{Llíria, España}{abr 2023 – may 2026}
\begin{itemize}
  \item Reduje el consumo de \textit{stack} de las tareas sobre RTOS entre un 30\% y un 50\%, optimizando el uso de memoria en los STM32.
  \item Optimicé el manejo de peticiones Modbus TCP, reduciendo el uso de CPU bajo carga (40.000 peticiones/s) del 80\% al 17\% sin pérdida de respuestas.
  \item Mejoré la estabilidad del sistema de archivos embebido, reduciendo los tiempos de acceso en operación continua.
  \item Implementé y mantuve comunicaciones industriales en C con capas de abstracción de hardware (HAL): Modbus RTU/TCP, UART, I2C y SPI.
  \item Automaticé miles de pruebas unitarias (Ceedling) y de integración sobre hardware real (pytest): rendimiento TCP/FTP y seguridad FTP.
  \item Revisé los \textit{merge requests} del equipo, detectando bugs y proponiendo mejoras antes de su integración.
\end{itemize}\vspace{2pt}

\cventry{Investigador predoctoral (acústica y ultrasonidos)}{UPV · CSIC · i3M / IIS La Fe}{Valencia, España}{jul 2021 – abr 2023}
\begin{itemize}
  \item Diseñé metamateriales acústicos (metadifusores y lentes) que reducen el tamaño de los difusores convencionales, validados con simulación en MATLAB y COMSOL.
  \item Fabriqué y caractericé prototipos; publiqué resultados en congresos internacionales (Euronoise, Tecniacústica).
  \item Tutoricé el Trabajo Fin de Máster de un estudiante del Máster en Ingeniería Acústica de la UPV, desde el diseño de la investigación hasta la defensa.
\end{itemize}\vspace{2pt}

\cventry{Investigador en proyecto europeo}{UPV · Agencia Espacial Europea (ESA) — IGIC}{Gandía, España}{feb 2020 – jul 2021}
\begin{itemize}
  \item Diseñé y validé experimentalmente un metamaterial de mitigación de ruido para el lanzamiento del cohete VEGA (proyecto ESA AO/1-9479/18/NL/LvH), cumpliendo el 100\% de los objetivos del proyecto y con resultados aceptados por la ESA.
  \item Desarrollé software de medición acústica según ISO 10534-2 y ASTM 2611 (\href{https://github.com/jmrplens/A-Lab}{A\,|\,Lab}) y un sistema robótico de 3 ejes para mediciones experimentales.
  \item Coautor de publicaciones/ponencias internacionales (Euronoise, ECSSMET).
\end{itemize}\vspace{2pt}

% =====================================================================
%  COMPETENCIAS TÉCNICAS
% =====================================================================
\section{Competencias Técnicas}
\cvskill{Firmware y embebido}{C, STM32/ARM, ESP32/ESP-IDF, FreeRTOS, PlatformIO, Modbus RTU/TCP, I2C/SPI/UART, JTAG/SWD, Secure Boot, Doxygen, instrumentación (osciloscopio, generador de señal, analizador lógico)}
\cvskill{Software y lenguajes}{Python, Go, C++, Bash, TypeScript, Git, Docker, CMake, Clang tooling}
\cvskill{Calidad, CI/CD y DevSecOps}{GitHub Actions, GitLab CI, Jenkins, SonarQube, GoReleaser, Ceedling, pytest, Playwright, análisis estático/SAST (CodeQL, gitleaks), pre-commit}
\cvskill{Redes}{TCP/IP, MikroTik RouterOS, WireGuard, IPv6, QUIC/HTTP3, DNS}
\cvskill{Seguridad y criptografía}{TLS/mTLS, PKI (infraestructura de clave pública), criptografía aplicada (AEAD, derivación de claves/KDF), hardening, CrowdSec}
\cvskill{IA y agentes}{Model Context Protocol (MCP), agentes de IA para programar (Claude Code, Copilot), integración de LLMs}
\cvskill{Simulación y medida}{MATLAB, COMSOL Multiphysics, LabVIEW, instrumentación acústica (acelerómetros láser y de vibración, micrófonos)}

% =====================================================================
%  PROYECTOS OPEN SOURCE
% =====================================================================
\section{Proyectos Open Source}
\cvproject{\href{https://github.com/jmrplens/gitlab-mcp-server}{gitlab-mcp-server}}{Go · servidor MCP (Model Context Protocol)}{14\,\faStar \ \divider\ \texttt{\textasciitilde}23k descargas \divider\ Glama / mcp.so / Cursor}
Servidor MCP de GitLab que expone más de 1.000 operaciones a asistentes de IA mediante un diseño dinámico de 2 herramientas (\textit{find/execute}), con transportes stdio/HTTP/OAuth.\par\vspace{4pt}

\cvproject{\href{https://github.com/jmrplens/PyOctaveBand}{PyOctaveBand}}{Python · procesado de señal}{96\,\faStar \ \divider\ PyPI}
Biblioteca de filtrado por bandas de octava y fracciones de octava para señales en el dominio del tiempo.\par\vspace{4pt}

\cvproject{\href{https://github.com/jmrplens/cs-routeros-bouncer}{cs-routeros-bouncer}}{Go · seguridad de red}{10\,\faStar}
Bouncer de CrowdSec para MikroTik RouterOS; gestiona \textit{address lists} y reglas de firewall vía la API de RouterOS.\par\vspace{4pt}

\cvproject{\href{https://github.com/jmrplens/Cloudflare-DNS-Updater}{Cloudflare-DNS-Updater}}{Bash · DevOps}{26\,\faStar}
Cliente de DNS dinámico para Cloudflare (IPv4/IPv6) con batería de pruebas (bats) y cobertura.

% =====================================================================
%  FORMACIÓN
% =====================================================================
\section{Formación}
\cventry{Programa de Doctorado en Tecnologías para la Salud y el Bienestar}{Universitat Politècnica de València (i3M / UMIL)}{Valencia, España}{2023}
{\color{muted}\small Investigación en metamateriales acústicos para imagen médica y control de haces ultrasónicos (programa no finalizado; reorientación a desarrollo de software industrial).}\par\vspace{4pt}

\cventry{Máster en Ingeniería Acústica}{Universitat Politècnica de València (EPSG)}{Gandía, España}{2019}
{\color{muted}\small Titulación con honores; TFM con mención especial.}\par\vspace{4pt}

\cventry{Grado en Ingeniería de Telecomunicación (Sonido e Imagen)}{Universidad de Alicante (EPS)}{Alicante, España}{2018}
{\color{muted}\small Dos especializaciones: Ingeniería Acústica y Tecnología Audiovisual.}\par\vspace{4pt}

\cventry{Téc. Superior en Sonido \& Téc. en Sistemas Electrónicos (CFGS/CFGM)}{IES Salesianos · IES Luis García Berlanga}{Elche / Alicante, España}{2009 – 2011}

\end{document}
